\section*{Hoofdstuk \ref{ch:g11:vision-and-brain} --- Zien en de hersenen}

\begin{solution}{exo:g11:vision-and-brain:1}
Licht uit het linkerveld valt op de rechterhelft van elk netvlies; de
vezels van die rechterhelft steken bij het linkeroog over bij de
kruising, die van het rechteroog blijven rechts; samen vormen zij de
rechter oogzenuwbaan, schakelen zij over in de thalamus en bereiken zij de
rechter \omterm{prop:g11:vision-and-brain:pathway}{primaire visuele schors}.
\end{solution}

\begin{solution}{exo:g11:vision-and-brain:2}
De twee oogzenuwen komen samen; de vezels van de neuszijdige helft van elk
netvlies steken over naar de andere kant, die van de slaapzijdige helft
niet. Daarvoorbij draagt elke baan één helft van het gezichtsveld.
\end{solution}

\begin{solution}{exo:g11:vision-and-brain:3}
Het kleurgebied (de wereld in grijs gezien), het bewegingsgebied (geen
waarneming van beweging), het gezichtengebied (gezichten worden niet
herkend).
\end{solution}

\begin{solution}{exo:g11:vision-and-brain:4}
Plasticiteit: het vermogen van verbindingen tussen neuronen om naar
gebruik te worden versterkt, verzwakt, gevormd of gesnoeid. Gevoelige
periode: de vroege levensfase waarin ervaring de schors beslissend en
onomkeerbaar vormt.
\end{solution}

\begin{solution}{exo:g11:vision-and-brain:5}
Het bindt zich aan de serotoninereceptoren van neuronen in de visuele
gebieden, verstoort de stroom van signalen en brengt hallucinaties voort
zonder dat er licht bij hoort.
\end{solution}

\begin{solution}{exo:g11:vision-and-brain:6}
Na de kruising, links: de linker oogzenuwbaan, de linker overschakeling in
de thalamus of de linker \omterm{prop:g11:vision-and-brain:pathway}{primaire visuele schors}.
\end{solution}

\begin{solution}{exo:g11:vision-and-brain:7}
De overstekende vezels komen uit de neuszijdige helften van beide
netvliezen, die de buitenste (slaapzijdige) helften van het veld zien: de
patiënt verliest de linkerhelft van het veld van het linkeroog en de
rechterhelft van dat van het rechteroog --- kokerzicht aan beide zijden.
\end{solution}

\begin{solution}{exo:g11:vision-and-brain:8}
Normaal: $12 + 20 + 36 + 20 = 88\%$. Beroofd: $2 + 3 + 5 + 30 = 40\%$,
waarvan slechts 10\% overwegend.
\end{solution}

\begin{solution}{exo:g11:vision-and-brain:9}
Afplakken dwingt de schors de invoer van het zwakke oog te gebruiken,
zodat zijn verbindingen tijdens de gevoelige periode behouden en versterkt
worden in plaats van gesnoeid. Bij een volwassene is die periode voorbij:
verbindingen worden op die schaal niet meer omgelegd, dus kan het zwakke
oog zijn gebied niet terugkrijgen.
\end{solution}

\begin{solution}{exo:g11:vision-and-brain:10}
Haar \omterm{def:g11:the-eye:photoreceptors}{kegeltjes} en hun \omterm{def:g11:the-eye:photoreceptors}{opsines} zijn normaal; wat verloren is, is de fase in
de schors die hun verhoudingen uitlegt. Een elektrische meting van de
reactie van het netvlies op gekleurde lichtbronnen zou bij haar normaal
zijn en bij een kleurenblinde afwijkend; en haar uitval is volledig en
verworven, die van hen gedeeltelijk en aangeboren.
\end{solution}

\begin{solution}{exo:g11:vision-and-brain:11}
Jarenlang oefenen heeft de verbindingen voor haar vingers versterkt en
vermenigvuldigd, waardoor hun gebied in de schors groter is geworden.
Zonder oefenen zou dat gebied over jaren weer krimpen, zoals de
jongleerstudies laten zien.
\end{solution}

\begin{solution}{exo:g11:vision-and-brain:12}
In het begin ziet zij licht, kleuren en bewegende vlekken maar herkent zij
geen vormen, gezichten of diepte: het netvlies werkt, maar de schors heeft
tijdens de gevoelige periode nooit de verbindingen gebouwd die zijn
signalen uitleggen. Zij zal leren het zien gedeeltelijk en langzaam te
gebruiken, maar nooit met het gemak van iemand die vanaf de geboorte zag.
\end{solution}

\begin{solution}{exo:g11:vision-and-brain:13}
In de gele vlek heeft elk \omterm{def:g11:the-eye:photoreceptors}{kegeltje} zijn eigen ganglioncel, dus stuurt het
midden van het veld per graad veel meer vezels dan de rand, waar honderd
\omterm{def:g11:the-eye:photoreceptors}{staafjes} er één delen; de schors verdeelt haar oppervlak naar het aantal
vezels, en dus naar detail, niet naar graden van het veld.
\end{solution}

\begin{solution}{exo:g11:vision-and-brain:14}
De werking van dat gebied ís de waarneming van kleur, of het signaal nu
van het oog of uit het geheugen komt; de waarneming is wat de gebieden
doen, niet wat het oog aanlevert. Een hallucinatie is diezelfde werking,
door een middel in gang gezet in plaats van door de wil of door licht.
\end{solution}

\begin{solution}{exo:g11:vision-and-brain:15}
Het oog zet licht om in signalen; het zien wordt door de hersenen gedaan.
Een patiënt met gave ogen kan door een beschadiging in de schors kleur,
beweging of gezichten kwijtraken; een katje met gave ogen wordt aan één
oog in de praktijk blind als de schors het niet mag aansluiten; en een
middel dat noch het oog noch de bouw van de schors raakt, laat iemand zien
wat er niet is. Het oog levert; de hersenen zien.
\end{solution}

\begin{solution}{pb:g11:vision-and-brain:1}
\textbf{1.} Het linkeroog of zijn oogzenuw, vóór de kruising: er zijn
alleen vezels van één oog getroffen.

\textbf{2.} De rechter oogzenuwbaan, thalamus of \omterm{prop:g11:vision-and-brain:pathway}{visuele schors}: een
uitval die beide ogen voor dezelfde veldhelft delen, kan alleen liggen
waar de vezels van de twee ogen voor die helft zijn samengebracht, dus na
de kruising.

\textbf{3.} De kruising zelf: de overstekende vezels, uit de neuszijdige
helften van beide netvliezen, die de slaapzijdige helften van het veld
zien, zijn doorgesneden.

\textbf{4.} Het bewegingsgebied, voorbij de primaire schors. Gave velden
en een gave gezichtsscherpte laten zien dat de primaire schors en alles
ervóór werken.

\textbf{5.} B, C en D hebben een gaaf netvlies; dat van A kan gaaf zijn of
niet. Ga dat na door het netvlies met een oogspiegel te bekijken en zijn
elektrische reactie op flitsen te meten.

\textbf{6.} De rechterschors ontvangt de linkerhelft van het veld.

\textbf{7.} Centrale $2^\circ$: \qty{1}{cm}; een band van $2^\circ$ op
$40^\circ$: \qty{0.1}{cm}. Verhouding 10.

\textbf{8.} De gele vlek pakt haar \omterm{def:g11:the-eye:photoreceptors}{kegeltjes} dicht opeen, elk met een
eigen ganglioncel, dus sturen de centrale graden veel meer vezels dan de
rand; de schors geeft hun naar verhouding meer oppervlak.

\textbf{9.} Een kleine blinde vlek in het midden van het linkerveld, wat
verlammend is voor het lezen; verder naar voren een groter blind gebied
aan de rand van het linkerveld, dat minder opvalt.

\textbf{10.} Door zijn ogen te bewegen brengt de patiënt elk woord op een
deel van het netvlies waarvan de schors gaaf is, en gebruikt hij de rand
van de beschadiging; traag, maar mogelijk.

\textbf{11.} $20 + 36 + 20 = 76\%$.

\textbf{12.} $2 + 3 + 5 + 30 = 40\%$, en slechts 10\% met het linkeroog
als voornaamste of gelijkwaardige invoer.

\textbf{13.} Zijn invoer bereikt de schors wel, maar zijn verbindingen
zijn bij gebrek aan gebruik verzwakt en gesnoeid; de verbindingen van het
rechteroog hebben zich in de vrijgekomen ruimte uitgebreid.

\textbf{14.} Dat de omlegging tot een gevoelige periode vroeg in het leven
beperkt blijft; de menselijke tegenhanger is amblyopie, die alleen in de
kindertijd ontstaat en daarna niet meer te behandelen is.

\textbf{15.} Doordat geen van beide ogen wordt bevoordeeld, is er geen
overname: het histogram blijft ruwweg symmetrisch maar met minder
reagerende neuronen in het geheel, en het zicht van het katje is aan beide
ogen slecht --- beroving zonder concurrentie verzwakt beide.

\textbf{16.} De hersenen onderdrukken het beeld van het gedraaide oog om
dubbelzien te vermijden; ongebruikt worden zijn verbindingen tijdens de
gevoelige periode gesnoeid en wordt het oog blijvend zwak. Een bril
verbetert de optiek, maar het oog wordt niet gebruikt, dus gaat het
snoeien door.

\textbf{17.} Het goede oog afdekken dwingt de schors de signalen van het
zwakke oog te gebruiken, waardoor zijn verbindingen behouden en versterkt
worden.

\textbf{18.} De gevoelige periode eindigt rond het zevende jaar; daarna
verdeelt de schors haar gebied niet meer op die schaal opnieuw, dus zijn
de verloren verbindingen van het zwakke oog niet meer op te bouwen.

\textbf{19.} Bij de volwassene liggen de verbindingen al vast en zijn zij
stabiel; de overname door het andere oog is een verschijnsel van de zich
ontwikkelende schors, tijdens de gevoelige periode, en niet van de rijpe.

\textbf{20.} A: het oog of de oogzenuw; C: de kruising; B: de rechterbaan,
de thalamus of de primaire schors; D: het bewegingsgebied daarvoorbij. De
plasticiteit van de schors, tot een gevoelige periode beperkt, laat het
afplakken bij het kind op zijn vierde werken en op zijn vijftiende falen.
\end{solution}
